% Lösungen Einheit 4 von 4 – Winkel von Geraden gegen Ebenen und Bogenmaße (zusammenbau v0.1)
\einheitenkopf{Einheit 4 von 4 · Winkel von Geraden gegen Ebenen und Bogenmaße}

\erg{16}{a) Sinus \quad b) $\lvert\vec u\rvert = 3$; $\mathrm{sin}\,\varphi = \frac{2}{3}$, $\varphi \approx 41{,}8^\circ$ \quad c) $\vec u = \overrightarrow{PS} = (-3 | -3 | 3)$, $\lvert\vec u\rvert = \sqrt{27}$; $\mathrm{sin}\,\varphi = \frac{3}{\sqrt{27}}$, $\varphi \approx 35{,}3^\circ$ \quad d) $\overrightarrow{AB} = (3 | 4 | 5)$, Projektion $(3 | 4 | 0)$; $\mathrm{cos}\,\varphi = \frac{25}{\sqrt{50} \cdot 5}$, $\varphi \approx 45{,}0^\circ$ \quad e) $r = 53 - 3 = 50$; halber Mittelpunktswinkel: $\mathrm{tan}\,\frac{\alpha}{2} = \frac{30}{40}$, $\alpha \approx 73{,}7^\circ$; Bogen $b = \frac{\alpha}{360^\circ} \cdot 2\pi \cdot 50 \approx 64{,}4$ m}
\erg{17}{a) Kante FS: $\mathrm{tan}\,\varphi_K = \frac{6}{3\sqrt{2}}$, $\varphi_K \approx 54{,}7^\circ$; Licht: $\mathrm{sin}\,\varphi_L = \frac{1}{\sqrt{3}}$, $\varphi_L \approx 35{,}3^\circ$; die Kante liegt in derselben senkrechten Ebene wie der Lichtstrahl durch S und ist steiler, also fällt der Schatten der Spitze neben die Pyramide}
\erg{18}{a) Fehler: der Kosinus liefert den Winkel zur Normalen; zur Ebene gehört der Sinus: $\mathrm{sin}\,\varphi = \frac{8}{9}$, $\varphi \approx 62{,}7^\circ$ \quad b) Richtungs- und Normalenvektor schließen den Winkel zur Normalen ein; der gesuchte Winkel zur Ebene ist dessen Komplement, und der Kosinus eines Winkels ist der Sinus seines Komplements \quad c) $\lvert\vec u\rvert = \sqrt{401}$; $\mathrm{sin}\,\varphi = \frac{1}{\sqrt{401}}$, $\varphi \approx 2{,}9^\circ$, also ja}
